\documentclass[conference]{IEEEtran}
\IEEEoverridecommandlockouts
\usepackage{cite}
\usepackage{amsmath,amssymb,amsfonts}
\usepackage{algorithmic}
\usepackage{graphicx}
\usepackage{textcomp}
\usepackage{xcolor}
\usepackage[hidelinks]{hyperref}

\def\BibTeX{{\rm B\kern-.05em{\sc i\kern-.025em b}\kern-.08em
    T\kern-.1667em\lower.7ex\hbox{E}\kern-.125emX}}

\begin{document}

\title{Optimizing Bidirectional LSTM for Financial Sentiment Analysis via Cross-Domain Transfer Learning and Structural Regularization}

\author{
    \IEEEauthorblockN{Marchesi Luca}
    \IEEEauthorblockA{\textit{Fondamenti dell'intelligenza artificiale} \\
        Parma, Italy \\
        luca.marchesi3@studenti.unipr.it}
    \and
    \IEEEauthorblockN{Minei Francesco}
    \IEEEauthorblockA{\textit{Fondamenti dell'intelligenza artificiale} \\
        Parma, Italy \\
        francesco.minei@studenti.unipr.it}
}

\maketitle

\begin{abstract}
    This paper investigates the optimization of a Bidirectional Long Short-Term Memory (Bi-LSTM) network for sentiment analysis in the financial domain. Leveraging cross-domain transfer learning from the IMDB movie review dataset to the Financial PhraseBank, we address critical challenges such as sequence padding drift and catastrophic forgetting. Our methodology introduces a Masked Mean Pooling architecture combined with discriminative fine-tuning and selective layer freezing. Experimental results demonstrate that a significant reduction in model capacity (90\% fewer parameters) leads to superior generalization, achieving a validation loss of 0.775 and effectively mitigating the overfitting observed in larger architectures.
\end{abstract}

\begin{IEEEkeywords}
    Sentiment Analysis, Bi-LSTM, Transfer Learning, Financial NLP, Regularization, Masked Pooling.
\end{IEEEkeywords}

\section{Introduction}
Sentiment analysis in the financial sector provides critical insights for market prediction and decision-making. However, the scarcity of labeled financial data often necessitates the use of Transfer Learning. In this study, we explore the adaptation of a Bi-LSTM model, pre-trained on the large-scale IMDB dataset, to the specialized Financial PhraseBank dataset. We identify and resolve structural inefficiencies, specifically focusing on the degradation of hidden states due to padding and the propensity of complex models to memorize small-scale training sets.

\section{Proposed Methodology}
Our approach focuses on structural and training-level optimizations to adapt a general-purpose sentiment model to specialized financial texts.

\subsection{Architectural Refinement: Masked Mean Pooling}
Financial sentences are significantly shorter than standard movie reviews, leading to excessive zero-padding. Standard LSTM implementations often suffer from state drift when extracting the last hidden state. We implemented a \textit{Masked Mean Pooling} layer that filters out padding tokens during feature aggregation:
\begin{equation}
    \text{Features} = \frac{\sum_{t=1}^{L} (h_t \cdot m_t)}{\sum_{t=1}^{L} m_t}
\end{equation}
where $h_t$ is the LSTM output at time $t$ and $m_t$ is a binary mask representing non-padding tokens. This ensures the classifier receives a signal uncontaminated by padding noise.

\subsection{Selective Transfer Learning}
To preserve the linguistic representations learned during pre-training on the IMDB dataset, we employed a selective freezing strategy. During the full fine-tuning phase, the Embedding layer remained frozen, while the LSTM and the MLP head were updated using discriminative learning rates ($10^{-4}$ for LSTM and $10^{-3}$ for MLP head).

\subsection{Optimization: AdamW and Weight Decay}
We transitioned from standard Adam to AdamW to decouple weight decay from the gradient update. This modification provided more stable regularization, particularly crucial when fine-tuning on a dataset of approximately 4,000 samples where the risk of overfitting is paramount.

\section{Experimental Analysis}
The model was evaluated on the Financial PhraseBank (50\% agreement). Initial tests with a large architecture (2 layers, 256 hidden units) showed extreme overfitting, with training accuracy reaching 99\% while validation accuracy plateaued at 70\%.

By applying \textit{Model Squeezing} (reducing to 1 layer and 128 hidden units), we observed that a smaller hypothesis space forced the network to learn more robust features. As shown in Table \ref{tab:results}, the simplified model achieved a lower validation loss (0.775), proving that for small-scale financial datasets, architectural simplicity is paramount for generalization.

\begin{table}[htbp]
    \caption{Performance and Complexity Comparison}
    \begin{center}
        \begin{tabular}{|l|c|c|c|}
            \hline
            \textbf{Architecture}    & \textbf{Parameters} & \textbf{Val. Acc} & \textbf{Val. Loss} \\
            \hline
            Bi-LSTM (256x2)          & 2,442,243           & 69.51\%           & 0.784              \\
            \textbf{Bi-LSTM (128x1)} & \textbf{268,803}    & \textbf{71.50\%}  & \textbf{0.775}     \\
            \hline
        \end{tabular}
        \label{tab:results}
    \end{center}
\end{table}

\section{Conclusion}
In this study, we demonstrated that optimizing a Bi-LSTM for financial sentiment analysis requires a careful balance between capacity and regularization. Through the implementation of Masked Pooling and AdamW-based weight decay, we successfully stabilized the transfer learning process. Our findings suggest that the 71-72\% accuracy range represents a performance ceiling for recurrent architectures on this dataset; further improvements would likely require attention-based architectures or large-scale language models like FinBERT.

\begin{thebibliography}{00}
    \bibitem{b1} S. Hochreiter and J. Schmidhuber, ``Long short-term memory,'' Neural computation, vol. 9, no. 8, pp. 1735--1780, 1997.
    \bibitem{b2} P. Malo, A. Sinha, P. Korhonen, J. Wallenius, and P. Takala, ``Good debt or bad debt: Detecting semantic orientations in economic texts,'' Journal of the Association for Information Science and Technology, vol. 65, no. 4, pp. 782--796, 2014.
    \bibitem{b3} I. Loshchilov and F. Hutter, ``Decoupled weight decay regularization,'' in Proc. Int. Conf. Learn. Represent. (ICLR), 2019.
    \bibitem{b4} A. L. Maas et al., ``Learning word vectors for sentiment analysis,'' in Proc. 49th Annu. Meeting Assoc. Comput. Linguistics, 2011, pp. 142--150.
\end{thebibliography}

\end{document}
